\chapter{Łańcuchy Markowa w czasie ciągłym}
\label{ch:ctmc}
\paragraph{Odsyłacze międzytomowe.} Łańcuchy w czasie dyskretnym opisuje tom I, rozdział~\ref{I-ch:markov}, a ich wykorzystanie w MCMC tom IV, rozdział~\ref{IV-ch:mcmc}.


\section{Dlaczego macierz przejścia nie wystarcza?}
W łańcuchu Markowa z czasu dyskretnego macierz $P=(p_{ij})$
opisuje jeden krok. W czasie ciągłym pojawiają się dwie odrębne
losowości: \emph{kiedy} następuje zmiana stanu oraz \emph{dokąd}
proces przechodzi. W modelu awarii i naprawy urządzenie może długo
pozostawać sprawne, a po awarii przebywać w stanie uszkodzenia
przez losowy czas. Sam ciąg odwiedzanych stanów nie opisuje
czasów trwania poszczególnych faz.

\begin{definicja}[Macierz intensywności]
Niech $E$ będzie skończonym lub przeliczalnym zbiorem stanów.
\emph{Konserwatywną macierzą intensywności} (macierzą $Q$)
nazywamy macierz $Q=(q_{ij})_{i,j\in E}$ spełniającą:
$q_{ij}\ge0$ dla $i\ne j$, $q_i:=-q_{ii}\in[0,\infty)$
oraz $\sum_{j\ne i}q_{ij}=q_i$ dla każdego $i\in E$.
Gdy $q_i>0$, definiujemy prawdopodobieństwa skoku
$r_{ij}=q_{ij}/q_i$ dla $j\ne i$ i $r_{ii}=0$.
Stan z $q_i=0$ jest pochłaniający; dla wygody
w macierzy łańcucha skoków przyjmujemy wówczas
$r_{ii}=1$ i $r_{ij}=0$ dla $j\ne i$.
\end{definicja}

Wyraz $q_{ij}$, $i\ne j$, jest natężeniem \emph{bezpośredniego}
przejścia z $i$ do $j$, a $q_i$ całkowitą intensywnością
opuszczenia stanu $i$. Zapis $q_{ii}=-q_i$ nie oznacza
ujemnego prawdopodobieństwa: elementy macierzy $Q$ są
intensywnościami, a nie prawdopodobieństwami.

\begin{przyklad}[Układ sprawny--uszkodzony]
Niech $E=\{0,1\}$, gdzie $0$ oznacza urządzenie sprawne,
a $1$ uszkodzone. Jeśli awarie występują z intensywnością
$\alpha>0$, a zakończenia napraw z intensywnością $\beta>0$,
to
\[
Q=\begin{pmatrix}-\alpha&\alpha\\\beta&-\beta\end{pmatrix}.
\]
Czas do awarii, rozpoczynając od stanu $0$, ma rozkład
$\operatorname{Exp}(\alpha)$; czas naprawy w stanie $1$
ma rozkład $\operatorname{Exp}(\beta)$. Po każdym przejściu
zaczyna się nowy, niezależny czas oczekiwania.
\end{przyklad}

\begin{center}
\begin{tikzpicture}[>=Latex,thick,node distance=2.4cm,
 st/.style={draw,circle,minimum size=11mm,fill=blue!5}]
\node[st] (zero) {$0$};
\node[st,right=of zero] (one) {$1$};
\node[st,right=of one] (two) {$2$};
\draw[->,bend left=22] (zero) to node[above] {$q_{01}=2$} (one);
\draw[->,bend left=22] (one) to node[below] {$q_{10}=1$} (zero);
\draw[->] (one) -- node[above] {$q_{12}=3$} (two);
\node[below=9mm of two] {pochłanianie};
\end{tikzpicture}
\par\smallskip\small Schemat: intensywności skoków dla trzech stanów.
W stanie $1$ łączna intensywność wynosi $1+3=4$.
\end{center}


\begin{uwaga}[Czym macierz intensywności różni się od macierzy przejścia?]
Jeżeli proces zaczyna w stanie $i$, to dla $j\ne i$
$\Prob_i(X_h=j)=q_{ij}h+O(h^2)$ przy $h\downarrow0$
\emph{w skończonym modelu}. Tymczasem dla dowolnego
ustalonego $t>0$ prawdopodobieństwo $p_{ij}(t)$ może zawierać
wkład wielu następujących po sobie zmian stanu.
Zatem $q_{ij}$ nie jest ani prawdopodobieństwem pojedynczego
skoku (nim jest $q_{ij}/q_i$), ani prawdopodobieństwem
znalezienia się w $j$ po czasie $t$ (nim jest $p_{ij}(t)$).
\end{uwaga}

\section{Konstrukcja przez łańcuch skoków i czasy przebywania}
\begin{twierdzenie}[Konstrukcja minimalnego procesu skokowego]
\label{thm:jump-construction}
Niech $E$ będzie zbiorem przeliczalnym, $Q$ konserwatywną
macierzą intensywności z $q_i<\infty$ dla każdego $i\in E$,
a $\mu$ rozkładem probabilistycznym na $E$.
Istnieje przestrzeń probabilistyczna, na której można
zdefiniować ciąg stanów $Y_0,Y_1,\ldots$, ciąg czasów skoków
$0=J_0<J_1<\cdots$ aż do ewentualnego pochłonięcia oraz
proces $X$ następująco:
\begin{itemize}
\item $Y_0\sim\mu$, a ze stanu $i$ o $q_i>0$ następny stan
ma rozkład $(r_{ij})_{j\in E}$;
\item przy wejściu do $i$ o $q_i>0$ losujemy czas przebywania
o rozkładzie $\operatorname{Exp}(q_i)$, warunkowo względem
bieżącego stanu niezależny od przeszłości i od wyboru
następnego stanu;
\item dla $J_n\le t<J_{n+1}$ ustalamy $X_t=Y_n$,
a po wejściu do stanu z $q_i=0$ proces pozostaje w nim na stałe.
\end{itemize}
Jeśli suma wszystkich kolejnych czasów przebywania jest skończona,
jej granicę $\zeta$ nazywamy \emph{czasem eksplozji}; dla $t\ge\zeta$
ustalamy $X_t=\partial$ w dodatkowym stanie cmentarnym
$\partial\notin E$. Tak otrzymany proces nazywamy
\emph{minimalnym procesem skokowym o macierzy $Q$}.
Na przedziale $[0,\zeta)$ jego trajektorie są schodkowe i
c\`adl\`ag, a czasy rzeczywistych skoków nie mają atomów
w skończonych chwilach.
\end{twierdzenie}
\begin{proof}
Dla każdego stanu $i$ z $q_i>0$ wiersz $(r_{ij})$ jest
rozkładem probabilistycznym. Używając niezależnych zmiennych
jednostajnych na $(0,1)$ oraz ustalonego wyliczenia stanów,
możemy metodą odwrotnej dystrybuanty wylosować następny stan
z tego rozkładu. Pozwala to rekurencyjnie skonstruować
$Y_0,Y_1,\ldots$; po pochłonięciu pozostawiamy
$Y_n$ w stanie pochłaniającym. Zmienne pomocnicze
możemy zrealizować na przeliczalnym iloczynie przestrzeni
probabilistycznych wraz z niezależnym stanem początkowym.

Niezależnie od całego łańcucha $Y$ losujemy ciąg
niezależnych zmiennych $Z_0,Z_1,\ldots$ o rozkładzie
$\operatorname{Exp}(1)$, a czas przebywania przy wizycie
w stanie $Y_n=i$ definiujemy jako $H_n=Z_n/q_i$,
jeśli $q_i>0$, oraz $H_n=\infty$, jeśli $q_i=0$.
Z równości
$\Prob(Z_n/q_i>t)=e^{-q_i t}$ wynika żądany rozkład
warunkowy. Ustalamy $J_0=0$ i $J_{n+1}=J_n+H_n$ w rozszerzonej
półprostej. W szczególności po pochłonięciu kolejne
czasy skoków są nieskończone. Definiujemy
$\zeta=\lim_n J_n$. Wszystkie $J_n$ i $\zeta$ są mierzalne.
Dla $i\in E$ i skończonego $t$ mamy
\[
\{X_t=i\}=\bigcup_{n\ge0}
\{Y_n=i,\ J_n\le t<J_{n+1}\},
\qquad \{X_t=\partial\}=\{\zeta\le t\}.
\]
Zatem $X_t$ jest zmienną losową na dyskretnej
przestrzeni $E\cup\{\partial\}$.

Jeśli $J_n\uparrow\zeta<\infty$, to przed chwilą $\zeta$
występuje nieskończenie wiele skoków. Nie należy wtedy
arbitralnie przedłużać trajektorii w $E$, bo może ona nie
mieć granicy. Dodanie stanu $\partial$ daje jednoznaczny
minimalny proces do końca czasu. Każdy skończony $H_n>0$, więc
kolejne skończone chwile $J_n$ są ściśle rosnące; na $[0,\zeta)$
na każdym przedziale domkniętym zawartym w $[0,\zeta)$
występuje tylko skończenie wiele skoków. Stąd wskazana
regularność trajektorii przed eksplozją.
Dla $n\ge1$ i ustalonego skończonego $t$ mamy
$\Prob(J_n=t)=0$: warunkowo względem ciągu stanów
$Y$, jeśli $J_n<\infty$, jest on sumą $n$ niezależnych
zmiennych wykładniczych o dodatnich parametrach, a więc
ma gęstość (przez splot gęstości). Możliwy jest natomiast
atom w nieskończoności.
\end{proof}

\begin{uwaga}[Dwie różne macierze]
Macierz $R=(r_{ij})$ opisuje \emph{dokąd} idziemy po skoku.
Macierz $Q$ opisuje łącznie wybór nowego stanu i tempo
opuszczania poprzedniego. Nawet jeśli dwa procesy mają taki sam
łańcuch skoków $Y_n$, lecz różne intensywności $q_i$,
mają inne rozkłady $X_t$ w czasie kalendarzowym.
\end{uwaga}


\begin{przyklad}[Dlaczego rozkład wykładniczy jest naturalny dla jednorodnego łańcucha?]
Przypuśćmy, że układ zaczyna w stanie $0$, czeka losowy
czas $H$, a następnie nieodwracalnie przechodzi do $1$.
Jeżeli $H$ ma rozkład jednostajny na $[0,2]$, to przy $0<s<s+h<2$
\[
\Prob(H>s+h\mid H>s)=\frac{2-s-h}{2-s}.
\]
To prawdopodobieństwo zależy od czasu $s$, który upłynął
od uruchomienia. Dla procesu, który \emph{jednorodnie}
zachowuje się w stanie $0$, oczekiwalibyśmy natomiast,
aby przy obecnym stanie $0$ szansa pozostania w nim przez
kolejne $h$ jednostek czasu nie zależała od wieku pobytu
w tym stanie. Gdy $H\sim\operatorname{Exp}(q)$, rzeczywiście
\[
\Prob(H>s+h\mid H>s)=e^{-qh}.
\]
Równanie braku pamięci wyjaśnia, dlaczego w konstrukcji
jednorodnego łańcucha Markowa wybieramy czasy wykładnicze.
Przy innych prawach czasu przebywania można otrzymać
\emph{proces pół-Markowa}, w którym do stanu należy
dodać informację o czasie od ostatniej zmiany.
\end{przyklad}

\section{Eksplozja i jej wykluczanie}
\begin{twierdzenie}[Ograniczone intensywności wykluczają eksplozję]
\label{thm:bounded-rates}
Niech proces minimalny z twierdzenia~\ref{thm:jump-construction}
ma $\sup_{i\in E}q_i\le\Lambda<\infty$.
Wtedy $\Prob(\zeta=\infty)=1$.
\end{twierdzenie}
\begin{proof}
Jeśli $\Lambda=0$, wszystkie stany są pochłaniające
i teza jest natychmiastowa. Załóżmy $\Lambda>0$.
Jeżeli proces wchodzi do stanu pochłaniającego, to
nie może mieć nieskończenie wielu rzeczywistych skoków
w czasie skończonym. W przeciwnym razie kolejne czasy
przebywania są równe $H_n=Z_n/q_{Y_n}$, gdzie
$Z_n$ są niezależne o rozkładzie $\operatorname{Exp}(1)$.
Ponieważ $q_{Y_n}\le\Lambda$, mamy
$H_n\ge Z_n/\Lambda$ prawie na pewno.
Z lematu o braku eksplozji dla niezależnych wykładniczych
z rozdziału 4 wynika $\sum_n Z_n/\Lambda=\infty$
prawie na pewno. Także $\sum_n H_n=\infty$.
\end{proof}

\begin{przyklad}[Dlaczego skończoność każdego $q_i$ nie wystarcza?]
Niech $E=\N_0$, $X_0=0$, a jedyne dozwolone przejście ma postać
$n\to n+1$ z intensywnością $q_{n,n+1}=\lambda_n>0$.
Czas eksplozji jest sumą niezależnych zmiennych
$H_n\sim\operatorname{Exp}(\lambda_n)$.
Dla $\lambda_n=(n+1)^2$ mamy
\[
\E\zeta=\sum_{n=0}^\infty\frac1{(n+1)^2}<\infty.
\]
Ponieważ $\zeta\ge0$, skończona wartość oczekiwana
implikuje $\zeta<\infty$ prawie na pewno.
Nieskończenie wiele skoków występuje więc przed pewną
skończoną, losową chwilą! To model eksplozji.
Natomiast dla $\lambda_n=n+1$ zbieżność $\sum1/\lambda_n$
nie zachodzi; poniższe twierdzenie pokazuje brak eksplozji.
\end{przyklad}

\begin{twierdzenie}[Kryterium eksplozji procesu czystych urodzin]
\label{thm:pure-birth-explosion}
Niech $E=\N_0$, $X_0=0$, a jedyną możliwą zmianą stanu
$n$ będzie przejście do $n+1$ z intensywnością
$\lambda_n\in(0,\infty)$. Wówczas
\[
\Prob(\zeta<\infty)=1
\quad\Longleftrightarrow\quad
\sum_{n=0}^\infty\frac1{\lambda_n}<\infty.
\]
Jeśli szereg jest rozbieżny, to $\Prob(\zeta=\infty)=1$.
\end{twierdzenie}
\begin{proof}
Mamy $\zeta=\sum_{n\ge0}H_n$ z niezależnymi
$H_n\sim\operatorname{Exp}(\lambda_n)$.
Jeżeli $\sum1/\lambda_n<\infty$, to
$\E\zeta=\sum\E H_n<\infty$ na mocy twierdzenia o
zbieżności monotonicznej, skąd $\zeta<\infty$ prawie na pewno.

Załóżmy teraz $\sum1/\lambda_n=\infty$.
Dla $c>0$ niezależność i twierdzenie o zbieżności
zmajoryzowanej (przez stałą $1$) dają
\[
\E e^{-c\zeta}
=\lim_{m\to\infty}\prod_{n=0}^m\E e^{-cH_n}
=\prod_{n=0}^\infty\frac{\lambda_n}{\lambda_n+c}.
\]
Musimy pokazać, że iloczyn jest zerowy. Jego ujemny logarytm
to szereg $\sum_n\log(1+c/\lambda_n)$.
Jeśli nieskończenie wiele $\lambda_n\le c$, to w szeregu
nieskończenie wiele wyrazów jest nie mniejszych od $\log2$.
W przeciwnym razie dla wszystkich dostatecznie dużych $n$
mamy $0<c/\lambda_n<1$, a więc
$\log(1+c/\lambda_n)\ge c/(2\lambda_n)$.
W obu przypadkach szereg logarytmów jest rozbieżny.
Zatem $\E e^{-c\zeta}=0$.
Przy przyjętej konwencji $e^{-c\infty}=0$, a dla każdej
skończonej wartości $\zeta$ wykładnik jest ściśle dodatni.
Stąd $\Prob(\zeta<\infty)=0$, co kończy dowód.
\end{proof}

\section{Macierz przejścia i równania Kołmogorowa}
Teraz ograniczymy się do \emph{skończonej} przestrzeni $E$.
Założenie jest istotne: w nieskończonym zbiorze stanów
potrzebujemy dodatkowej kontroli, aby uzasadniać różniczkowanie
nieskończonych szeregów i wykluczać eksplozję.

\begin{definicja}[Macierz przejścia w czasie ciągłym]
Niech $E$ będzie skończone, a $X$ procesem o intensywnościach
$Q$. Dla $i,j\in E$ i $t\ge0$ definiujemy
$p_{ij}(t)=\Prob(X_t=j\mid X_0=i)$,
przy czym po prawej stronie rozumiemy prawo procesu
\emph{uruchomionego ze stanu $i$}, również wtedy,
gdy w wyjściowym rozkładzie początkowym $\Prob(X_0=i)=0$.
Macierz $P(t)=(p_{ij}(t))_{i,j\in E}$ nazywamy
\emph{macierzą przejścia w czasie $t$}.
\end{definicja}

\section{Zegar Poissona i uniformizacja}
Najwygodniejsza metoda konstrukcji skończonego łańcucha
pozwala zarazem udowodnić równania różniczkowe macierzy
przejścia bez odwoływania się do formalnej ``rachunkowości
prawdopodobieństw $o(h)$''.

\begin{twierdzenie}[Uniformizacja]
\label{thm:uniformization}
Niech $E$ będzie niepustym zbiorem skończonym,
a $Q$ konserwatywną macierzą intensywności.
Wybierzmy $\Lambda>\max_{i\in E}q_i$ i ustalmy
$R:=I+Q/\Lambda$. Wtedy $R$ jest macierzą stochastyczną.
Niech $(Y_n)_{n\ge0}$ będzie łańcuchem dyskretnym o macierzy
przejścia $R$, niezależnym od procesu Poissona $N$ o intensywności
$\Lambda$. Proces $X_t:=Y_{N_t}$ jest łańcuchem w czasie
ciągłym o intensywnościach $Q$, a jego macierz przejścia wynosi
\[
P(t)=e^{-\Lambda t}\sum_{n=0}^\infty
\frac{(\Lambda t)^n}{n!}R^n=e^{tQ},\qquad t\ge0.
\]
Dla każdego $i\ne j$ zachodzi
$\lim_{h\downarrow0}p_{ij}(h)/h=q_{ij}$,
a dla $i=j$ zachodzi
$\lim_{h\downarrow0}(p_{ii}(h)-1)/h=q_{ii}$.
\end{twierdzenie}
\begin{proof}
Dla $i\ne j$ mamy $R_{ij}=q_{ij}/\Lambda\ge0$,
a na przekątnej $R_{ii}=1-q_i/\Lambda>0$.
Suma każdego wiersza jest równa
$1+\Lambda^{-1}\sum_j q_{ij}=1$.
Wprowadzamy niezależny pomocniczy łańcuch dyskretny
i zegar Poissona.
W chwili zegara $Y$ dokonuje kroku macierzy $R$; może
przy tym pozostać w obecnym stanie. Te pozorne skoki
nazywamy \emph{skokami fikcyjnymi}.

Warunkowo względem $N_t=n$ końcowy stan procesu ma
rozkład $n$-krokowy $R^n$. Dlatego
$P(t)=\sum_n\Prob(N_t=n)R^n$, czyli pierwsza suma.
Ponieważ $I$ i $R$ komutują,
\[
e^{-\Lambda t}\sum_{n=0}^\infty\frac{(\Lambda tR)^n}{n!}
=e^{-\Lambda t I}e^{\Lambda tR}
=e^{\Lambda t(R-I)}=e^{tQ}.
\]
W podmultiplikatywnej normie macierzowej z $\norm{I}=1$
(np.\ normie operatorowej) szereg jest bezwzględnie zbieżny, gdyż
$\sum_{n\ge0}(\Lambda t)^n\norm{R^n}/n!$ można oszacować
przez $\exp(\Lambda t\norm{R})$, korzystając z
$\norm{R^n}\le\norm{R}^n$.
To uzasadnia również mnożenie szeregów wykładniczych.
Rozwinięcie wykładnika macierzowego
$e^{hQ}=I+hQ+O(h^2)$ przy $h\downarrow0$
daje oba wskazane limity.

Pozostaje porównać konstrukcję z rzeczywistymi
czasami przebywania. Przy starcie ze stanu $i$ o $q_i>0$
niech $K=\inf\{n\ge1:Y_n\ne i\}$.
Z własności Markowa łańcucha dyskretnego
\[
\Prob_i(K=k,Y_K=j)=R_{ii}^{k-1}R_{ij},
\qquad k\ge1,\ j\ne i.
\]
Stąd $K$ ma rozkład geometryczny z parametrem
$1-R_{ii}=q_i/\Lambda$, a kierunek $Y_K$ jest niezależny
od $K$ i ma rozkład $q_{ij}/q_i$.
Ponieważ cały łańcuch $Y$ jest niezależny od zegara,
czas pierwszego rzeczywistego skoku $H=T_K$, gdzie
$T_k$ są chwilami impulsów Poissona, jest niezależny
od kierunku $Y_K$. Ponadto
\[
\Prob_i(H>t)=\sum_{n\ge0}\Prob(N_t=n)R_{ii}^n
=e^{-\Lambda t(1-R_{ii})}=e^{-q_i t}.
\]
Jeżeli $q_i=0$, wiersz $R$ ma jedynkę na przekątnej
i żadna rzeczywista zmiana stanu nie następuje.
Po rzeczywistym skoku argument można powtórzyć:
łańcuch dyskretny ma mocną własność Markowa w indeksie
$K$ (tom I), a niezależne wykładnicze odstępy zegara
po tym indeksie są niezależne od dotychczasowych.
Ostatnia własność wynika przez rozbicie względem
$\{K=k\}$, ponieważ $K$ zależy wyłącznie od $Y$,
a odstępy zegara są niezależne od $Y$ i między sobą.
Iteracja daje warunkowe prawa kolejnych czasów i kierunków
takie jak w twierdzeniu~\ref{thm:jump-construction}.
Tym samym uniformizacja realizuje ten sam model.
\end{proof}


\begin{uwaga}[Dlaczego fikcyjne impulsy nie zmieniają modelu?]
Jeżeli proces zaczyna w stanie $i$ z $q_i>0$,
liczba impulsów fikcyjnych przed pierwszą rzeczywistą
zmianą stanu jest geometryczna. Można obliczyć
prawdopodobieństwo pozostawania w $i$ \emph{bez
rzeczywistego skoku} do chwili $t$ bez powoływania się
jedynie na intuicję pamięci rozkładu wykładniczego:
\[
\sum_{n=0}^\infty\Prob(N_t=n)R_{ii}^{n}
=e^{-\Lambda t}\exp(\Lambda tR_{ii})
=e^{-q_it}.
\]
Tak otrzymujemy czas oczekiwania $\operatorname{Exp}(q_i)$.
Warunkowo względem tego, że impuls powoduje rzeczywisty skok,
prawdopodobieństwo wybrania stanu $j\ne i$ jest
$R_{ij}/(1-R_{ii})=q_{ij}/q_i$.
Po skoku powtarzamy argument warunkowo względem
nowego stanu, korzystając z uzasadnienia w dowodzie
uniformizacji. Ta obserwacja uzasadnia
równoważność uniformizacji i konstrukcji przez czasy przebywania.
\end{uwaga}

Ważna korzyść uniformizacji jest obliczeniowa: możemy
symulować jeden wspólny zegar Poissona i niezależną sekwencję
kroków macierzy $R$, a przy prawdopodobieństwach przejścia
stosować nieujemny szereg zamiast bezpośrednio obliczać
wykładnik macierzowy.

\begin{twierdzenie}[Równania Chapmana--Kołmogorowa i Kołmogorowa]
\label{thm:kolmogorov-eq}
Dla skończonej przestrzeni stanów i konserwatywnej macierzy
intensywności $Q$ macierz przejścia $P(t)=e^{tQ}$ spełnia
$P(0)=I$ i $P(t+s)=P(t)P(s)$ dla $s,t\ge0$. Ponadto
jest różniczkowalna dla każdego $t\ge0$ (w zerze prawostronnie)
i spełnia \emph{równanie wsteczne} oraz \emph{równanie w przód}:
\[
P'(t)=QP(t),\qquad P'(t)=P(t)Q.
\]
Współrzędnie, dla każdego $i,j\in E$,
\[
\frac{\dd}{\dd t}p_{ij}(t)=\sum_{k\in E}q_{ik}p_{kj}(t),
\qquad
\frac{\dd}{\dd t}p_{ij}(t)=\sum_{k\in E}p_{ik}(t)q_{kj}.
\]
\end{twierdzenie}
\begin{proof}
Wzór $P(t)=e^{tQ}$ otrzymaliśmy z konstrukcji, zatem
$P(t+s)=e^{(t+s)Q}=e^{tQ}e^{sQ}$.
Istnieje też interpretacja probabilistyczna: warunkujemy
stan w chwili $t$ i sumujemy po wszystkich stanach
pośrednich $k$, otrzymując
$p_{ij}(t+s)=\sum_kp_{ik}(t)p_{kj}(s)$.
Na każdym przedziale $[0,T]$ zarówno szereg wykładnika
macierzowego, jak i szereg jego pochodnych są jednostajnie
zbieżne: w normie operatorowej ich wyrazy są ograniczone
odpowiednio przez $(T\norm Q)^n/n!$ i, dla $n\ge1$,
$\norm Q(T\norm Q)^{n-1}/(n-1)!$.
Można więc różniczkować wyraz po wyrazie; dostajemy
$P'(t)=Qe^{tQ}=e^{tQ}Q$.
Przemienność $Q$ z własnym wykładnikiem wynika również
bezpośrednio z definicji szeregu.
Zapisując mnożenie macierzy współrzędnie, uzyskujemy
obydwa wzory skalarne.
\end{proof}

\begin{uwaga}[Dlaczego są dwa równania?]
Równanie wsteczne analizuje \emph{pierwszą decyzję po starcie}:
pierwszy indeks $i$ należy do stanu startowego.
Równanie w przód analizuje \emph{dopływ do stanu końcowego}:
drugi indeks $j$ jest stanem w chwili $t$.
W skończonym jednorodnym modelu obydwa równania są poprawne
i prowadzą do tej samej macierzy. Dla modeli z nieskończenie
wieloma stanami albo z eksplozją przejścia do różniczkowania
w nieskończonych szeregach wymagają osobnego dowodu.
\end{uwaga}


\begin{przyklad}[Powrót do stanu początkowego to nie brak skoku]
W modelu $E=\{0,1\}$ z intensywnościami
$0\to1$ równą $\alpha>0$ i $1\to0$ równą $\beta>0$
czas do pierwszego \emph{rzeczywistego} skoku, przy starcie
z $0$, ma rozkład $\operatorname{Exp}(\alpha)$. Zatem
\[
\Prob_0(\text{brak skoku do }t)=e^{-\alpha t}.
\]
Natomiast $p_{00}(t)=\Prob_0(X_t=0)$ uwzględnia również
ścieżki z dwoma, czterema lub większą parzystą liczbą skoków.
W istocie, z rozwiązania równań w następnym podrozdziale,
\[
 p_{00}(t)=\frac{\beta+\alpha e^{-(\alpha+\beta)t}}
 {\alpha+\beta}>e^{-\alpha t}\qquad(t>0).
\]
Nierówność jest ścisła, bo z dodatnim prawdopodobieństwem
występują dwa skoki $0\to1\to0$ przed chwilą $t$.
Rozróżnienie między ``stan nie zmienił się do chwili $t$''
a ``w chwili $t$ obserwuję stan początkowy'' będzie ważne
przy modelowaniu trajektorii i przy zadaniach o czasach trafienia.
\end{przyklad}

\section{Warunkowanie, własność Markowa i czasy zatrzymania}
\begin{twierdzenie}[Własność Markowa w czasie ciągłym]
Niech $E$ będzie skończone, $X$ procesem o intensywnościach $Q$,
a $\calF_t^X=\sigma(X_s:0\le s\le t)$ filtracją naturalną.
Dla każdego $s,t\ge0$ i każdej funkcji $f:E\to\R$
\[
\E[f(X_{s+t})\mid\calF_s^X]=(P(t)f)(X_s),\qquad
(P(t)f)(i)=\sum_{j\in E}p_{ij}(t)f(j).
\]
W konsekwencji dla $j\in E$
$\Prob(X_{s+t}=j\mid\calF_s^X)=p_{X_sj}(t)$.
\end{twierdzenie}
\begin{proof}
Pracujemy na przestrzeni z konstrukcji uniformizacyjnej:
$X_t=Y_{N_t}$, gdzie łańcuch dyskretny $Y$ o macierzy $R$
jest niezależny od zegara Poissona $N$ o intensywności $\Lambda$.
Niech $\calH_s$ oznacza informację wygenerowaną przez
$(N_u)_{0\le u\le s}$ oraz przez zmienne
$V_j^s$, równe $Y_j$ na $\{j\le N_s\}$ i ustalonemu
dodatkowemu symbolowi poza $E$ w przeciwnym razie.
Dla $u\le s$ zmienna $X_u=Y_{N_u}$ jest
$\calH_s$-mierzalna, gdyż na $\{N_u=j\}$ jest równa
$V_j^s$. Zatem $\calF_s^X\subseteq\calH_s$.
Dla każdego $A\in\calH_s$ i $n\ge0$ zachodzi
$A\cap\{N_s=n\}\in\sigma(\calF_s^N,Y_0,\ldots,Y_n)$:
klasa zdarzeń o tej własności jest $\sigma$-ciałem
i zawiera wszystkie wskazane generatory.

Ustalmy $k\ge0$ i weźmy zdarzenie $\{N_s=n\}$.
Niezależność całego procesu $Y$ od $N$ i własność
Markowa dyskretnego łańcucha dają, dla każdego stanu $i$,
\[
\E[f(Y_{n+k})\mid Y_0,\ldots,Y_n]=(R^kf)(Y_n).
\]
Warunkowo względem $\calH_s$ liczba dodatkowych impulsów
$K=N_{s+t}-N_s$ ma rozkład $\operatorname{Poisson}(\Lambda t)$
i jest niezależna od całego łańcucha $Y$ oraz przeszłości
zegara łącznie. Dokładniej, dla $A\in\calH_s$,
na $\{N_s=n\}$ stosujemy niezależność $K$ od
$\sigma(\calF_s^N,Y_j:j\ge0)$ oraz dyskretną własność
Markowa, rozszerzoną na $\sigma(\calF_s^N,Y_0,\ldots,Y_n)$
dzięki niezależności $Y$ od zegara (argument $\pi$--$\lambda$).
Otrzymujemy
\[
\E[\chi_{A\cap\{N_s=n\}}\chi_{\{K=k\}}f(Y_{n+k})]
=\Prob(K=k)\E[\chi_{A\cap\{N_s=n\}}(R^kf)(Y_n)].
\]
Sumując po $n,k$ i stosując charakterystykę
warunkowej wartości oczekiwanej, dostajemy:
\[
\begin{aligned}
\E[f(X_{s+t})\mid\calH_s]
 &=\sum_{k=0}^\infty
 e^{-\Lambda t}\frac{(\Lambda t)^k}{k!}
 (R^kf)(Y_{N_s})\\
 &= (P(t)f)(X_s).
\end{aligned}
\]
Powyższy rachunek można najpierw przeprowadzić na zdarzeniach
$\{N_s=n\}$, a potem zsumować po $n\ge0$; jest to dozwolone,
ponieważ $f$ jest ograniczona, $E$ jest skończone,
a szereg wag Poissona sumuje się do $1$.
Wreszcie $\calF_s^X\subseteq\calH_s$ i prawa strona
jest $\calF_s^X$-mierzalna. Własność wieży
\emph{warunkowej wartości oczekiwanej} daje
\[
\E[f(X_{s+t})\mid\calF_s^X]
=\E[(P(t)f)(X_s)\mid\calF_s^X]
=(P(t)f)(X_s).
\]
Dla $f=\chi_{\{j\}}$ otrzymujemy zapis
w terminach prawdopodobieństw warunkowych.
\end{proof}


\begin{twierdzenie}[Własność Markowa procesu minimalnego na przeliczalnej przestrzeni stanów]
Niech $X$ będzie procesem minimalnym z
twierdzenia~\ref{thm:jump-construction}, z dodatkowym
stanem pochłaniającym $\partial$. Dla $i\in E\cup\{\partial\}$
oznaczmy przez $P_t f(i)$ wartość oczekiwaną $f(X_t)$
dla procesu uruchomionego ze stanu $i$.
Dla każdej ograniczonej funkcji $f:E\cup\{\partial\}\to\R$
oraz $s,t\ge0$ zachodzi
\[
\E[f(X_{s+t})\mid\calF_s^X]=(P_t f)(X_s)
\quad\text{prawie na pewno}.
\]
\end{twierdzenie}
\begin{proof}
Wybierzmy rosnący ciąg skończonych zbiorów $E_m$
o sumie $E$. Proces $X^{(m)}$ otrzymujemy, wysyłając
$X$ do $\partial$ przy pierwszym wyjściu poza $E_m$
(również w chwili $0$, jeśli $X_0\notin E_m$).
Do wyjścia korzystamy z tych samych stanów i czasów
przebywania co w procesie pierwotnym.
Na $E_m\cup\{\partial\}$ macierz intensywności zachowuje
$q_{ij}$ dla $i,j\in E_m$ i ma
$q_{i\partial}=\sum_{j\notin E_m}q_{ij}$,
a wiersz $\partial$ jest zerowy. Konstrukcja czasów
przebywania i kierunków skoku daje dokładnie ten
skończony model, więc stosujemy udowodnioną już
własność Markowa w przypadku skończonym.

Na wspólnym zbiorze pełnej miary, na którym
$\sum_n Z_n=\infty$, eksplozja wymusza opuszczenie
każdego skończonego zbioru przed chwilą $\zeta$:
gdyby wszystkie odwiedzane stany należały do
pewnego skończonego zbioru $F$, mielibyśmy
$H_n\ge Z_n/\max_{i\in F}q_i$, jeśli maksimum jest
dodatnie, a przy maksimum zero nie byłoby skoków.
To wyklucza skończoną sumę czasów przebywania.
Dla ustalonej chwili $v<\zeta$ proces odwiedza
do $v$ tylko skończenie wiele stanów, toteż
$X_v^{(m)}=X_v$ dla dostatecznie dużego $m$.
Jeśli $\zeta\le v$, wszystkie obcięte procesy
są w $\partial$ w chwili $v$ na wskazanym zbiorze
pełnej miary. W konsekwencji dla każdej ustalonej
chwili $v$ zachodzi $X_v^{(m)}\to X_v$ prawie na pewno,
w topologii dyskretnej. Ten sam argument przy
starcie z ustalonego $i$ i twierdzenie o zbieżności
zmajoryzowanej dają $P_t^{(m)}f(i)\to P_t f(i)$,
gdy $m$ jest dostatecznie duże, by obejmować $i$
(lub $i=\partial$).

Dla $0\le u_1<\cdots<u_l\le s$ oraz dowolnych
zbiorów $B_1,\ldots,B_l\subseteq E\cup\{\partial\}$
oznaczmy $C_m=\bigcap_{j=1}^l\{X_{u_j}^{(m)}\in B_j\}$
i analogicznie $C$ dla $X$.
Skończony model spełnia
\[
\E[\chi_{C_m} f(X_{s+t}^{(m)})]
=\E[\chi_{C_m}(P_t^{(m)}f)(X_s^{(m)})].
\]
Na wspólnym zbiorze pełnej miary dla wymienionych
chwil mamy od pewnego $m$ równość $X_s^{(m)}=X_s$
oraz $\chi_{C_m}=\chi_C$. Obie strony są ograniczone
przez $\norm{f}_\infty$ pod znakiem wartości oczekiwanej,
zatem przejście do granicy daje
$\E[\chi_C f(X_{s+t})]=\E[\chi_C(P_tf)(X_s)]$.
Zdarzenia cylindryczne $C$ tworzą $\pi$-układ
generujący $\calF_s^X$. Argument $\pi$--$\lambda$
(przy funkcji $f$ dowolnego znaku stosujemy równość
całek jako równość miar skończonych po dodaniu stałej)
rozszerza tę tożsamość na całe $\calF_s^X$.
Mierzalność $(P_tf)(X_s)$ kończy dowód.
\end{proof}

\begin{lemat}[Własność Markowa po prawostronnym uzupełnieniu filtracji]
\label{lem:markov-aug}
Niech $X$ będzie minimalnym procesem skokowym
na przeliczalnym $E$ (z dodatkowym stanem $\partial$
w razie eksplozji). Przyjmijmy topologię dyskretną
na przestrzeni stanów $E'=E\cup\{\partial\}$.
Proces ma prawostronnie ciągłe trajektorie:
przed eksplozją wynika to z konstrukcji, a od chwili
eksplozji trajektoria jest stała. Niech $\calF_t^0=\sigma(X_s:0\le s\le t)$
oraz
$\calF_t=\bigcap_{u>t}\sigma(\calF_u^0\cup\calN)$,
gdzie $\calN$ jest rodziną wszystkich podzbiorów zbiorów
zerowych z bazowego $\sigma$-ciała $\mathcal A$
przestrzeni probabilistycznej $(\Omega,\mathcal A,\Prob)$,
którą zakładamy zupełną.
Wówczas dla $t,h\ge0$ i każdej ograniczonej funkcji $f:E'\to\R$
\[
\E[f(X_{t+h})\mid\calF_t]=(P(h)f)(X_t)
\quad\text{prawie na pewno}.
\]
\end{lemat}
\begin{proof}
Ustalmy $t,h$ i $A\in\calF_t$. Dla każdego $n\ge1$
połóżmy $u_n=t+1/n$. Ponieważ $A\in\calF_t\subseteq\sigma(\calF_{u_n}^0\cup\calN)$,
po odrzuceniu nieistotnych zbiorów zerowych możemy
stosować własność Markowa względem $\calF_{u_n}^0$
do ograniczonej zmiennej $f(X_{u_n+h})$:
\[
\E[\chi_A f(X_{u_n+h})]
 =\E[\chi_A(P(h)f)(X_{u_n})].
\]
Uzasadnijmy przejście przez uzupełnienie. Każde zdarzenie
$A\in\sigma(\calF_{u_n}^0\cup\calN)$ różni się
zbiorem zerowym od pewnego $A_n\in\calF_{u_n}^0$;
wynika to z tego, że klasa takich zdarzeń tworzy
$\sigma$-ciało zawierające oba generatory. Zatem
poprzednią równość, początkowo prawdziwą dla $A_n$,
można zastosować także do $A$.
Prawostronna ciągłość daje $X_{u_n+h}\to X_{t+h}$
oraz $X_{u_n}\to X_t$ niemal na pewno.
W topologii dyskretnej zbieżność stanów oznacza ich
równość od pewnego indeksu. Funkcje $f$ i $P(h)f$
są ograniczone (druga przez $\norm{f}_\infty$),
więc twierdzenie o zbieżności zmajoryzowanej pozwala
przejść do granicy.
Otrzymana równość całek dla każdego $A\in\calF_t$
oraz $\calF_t$-mierzalność $(P(h)f)(X_t)$
stanowią charakterystykę warunkowej wartości oczekiwanej.
\end{proof}

\begin{twierdzenie}[Mocna własność Markowa]
\label{thm:ctmc-strong}
Niech $X$ będzie minimalnym procesem skokowym
o macierzy intensywności $Q$ na przeliczalnym $E$,
z dodatkowym stanem $\partial$ w razie eksplozji.
Niech $\mathbb F=(\calF_t)_{t\ge0}$ będzie zwykłym
uzupełnieniem filtracji naturalnej (jak w poprzednim
lemacie), a przestrzeń probabilistyczna niech będzie zupełna.
Jeżeli $\tau$ jest czasem zatrzymania skończonym prawie na
pewno, to dla każdego $t\ge0$ i ograniczonej funkcji
$f:E\cup\{\partial\}\to\R$
\[
\E[f(X_{\tau+t})\mid\calF_\tau]=(P(t)f)(X_\tau)
\quad\text{prawie na pewno}.
\]
\end{twierdzenie}
\begin{proof}
Na zdarzeniu $\{\tau=\infty\}$, które należy do $\calF_0$
dzięki zupełności, możemy zastąpić $\tau$ przez $0$.
Tak otrzymujemy skończony punktowo czas zatrzymania;
$\sigma$-ciała w tych dwóch chwilach są identyczne,
gdyż przecięcia w ich definicjach różnią się wyłącznie
na zdarzeniu zerowym należącym wraz z podzbiorami do $\calF_0$.
Wartości $X_\tau$ i $X_{\tau+t}$ na tym zdarzeniu
ustalamy dowolnie w $E$. Najpierw niech $\tau$ przyjmuje wartości w przeliczalnym
zbiorze $\{s_k:k\ge1\}\subset[0,\infty)$.
Dla $A\in\calF_\tau$ zdarzenie
$A\cap\{\tau=s_k\}$ należy do $\calF_{s_k}$
(ponieważ $\{\tau=s_k\}=\{\tau\le s_k\}\setminus
\{\tau<s_k\}$, gdzie $\{\tau<s_k\}$ jest sumą
$\{\tau\le q\}$ po nieujemnych wymiernych $q<s_k$;
dla $s_k=0$ jest puste).
Stosując własność Markowa w deterministycznej chwili $s_k$
\emph{dla uzupełnionej filtracji} (lemat~\ref{lem:markov-aug})
i sumując po $k$, uzyskujemy
\[
\E[\chi_A f(X_{\tau+t})]
=\E[\chi_A(P(t)f)(X_\tau)].
\]
Dla ogólnego czasu $\tau<\infty$ prawie na pewno niech
$\tau_m=2^{-m}\lceil2^m\tau\rceil$.
Jest to czas zatrzymania przyjmujący wartości w zbiorze
$2^{-m}\N_0$: dla $v\ge0$ zdarzenie $\{\tau_m\le v\}$
jest równe $\{\tau\le2^{-m}\lfloor2^mv\rfloor\}$.
Ponadto $\tau_m\downarrow\tau$
prawie na pewno. Mamy $\calF_\tau\subseteq\calF_{\tau_m}$,
więc dla każdego $A\in\calF_\tau$ poprzednia równość
zastosowana do $\tau_m$ daje
\[
\E[\chi_A f(X_{\tau_m+t})]
=\E[\chi_A(P(t)f)(X_{\tau_m})].
\]
Trajektorie procesu są prawostronnie ciągłe także
w stanie cmentarnym. Dlatego w topologii dyskretnej
$X_{\tau_m+t}\to X_{\tau+t}$ i $X_{\tau_m}\to X_\tau$
prawie na pewno. Ponieważ $f$ i $P(t)f$ są ograniczone,
możemy zastosować twierdzenie o zbieżności
zmajoryzowanej po obu stronach.
W efekcie identyczna równość zachodzi dla $\tau$ i wszystkich
$A\in\calF_\tau$. Proces $X$ jest progresywnie mierzalny
jako adaptowany proces prawostronnie ciągły
(twierdzenie~\ref{thm:prog}), zatem $X_\tau$ jest
$\calF_\tau$-mierzalna na mocy
twierdzenia~\ref{thm:sampled}. Otrzymujemy więc
charakterystykę wskazanej warunkowej wartości oczekiwanej.
\end{proof}

\begin{uwaga}[Prawo całej przyszłej trajektorii]
Mocna własność Markowa w powyższej postaci dotyczy
jednej przyszłej chwili. Dla $0<t_1<\cdots<t_m$
oraz ograniczonych funkcji $f_1,\ldots,f_m$ stosujemy
ją kolejno w czasach zatrzymania $\tau+t_{m-1},
\ldots,\tau+t_1,\tau$ (dodanie deterministycznej
stałej nieujemnej zachowuje własność czasu zatrzymania).
Własność wieży daje
\[
\E\left[\prod_{j=1}^m f_j(X_{\tau+t_j})
\mid\calF_\tau\right]
=
P(t_1)\bigl(f_1 P(t_2-t_1)
(f_2\cdots P(t_m-t_{m-1})f_m)\bigr)(X_\tau),
\]
gdzie zagnieżdżenie rozumiemy rekurencyjnie od
$f_m$ wstecz, a mnożenie funkcji jest punktowe.
Dla $m=1$ prawa strona to $(P(t_1)f_1)(X_\tau)$.
Wybór funkcji charakterystycznych i argument
$\pi$--$\lambda$ dla cylindrów rozszerzają to na
warunkowe prawo całego przesuniętego procesu
względem $\sigma$-ciała cylindrycznego: jest ono
prawem procesu startującego z $X_\tau$.
\end{uwaga}

\begin{przyklad}[Pierwsza naprawa jako chwila losowa]
W modelu sprawny--uszkodzony niech $\tau$ będzie chwilą
pierwszego powrotu do stanu sprawnego po awarii, przy starcie
ze stanu $0$. Jeżeli $\alpha,\beta>0$, to $\tau<\infty$
prawie na pewno. Jest to czas zatrzymania: dla
$t>0$ zdarzenie $\{\tau\le t\}$ jest sumą po
wymiernych $r\in(0,t)$ zdarzeń
$\{X_r=1\}\cap(\{X_t=0\}\cup
\bigcup_{q\in\mathbb Q\cap(r,t)}\{X_q=0\})$.
Równość wynika z prawostronnej ciągłości dyskretnych
trajektorii i dodatnich czasów przebywania; dla $t=0$
zdarzenie jest puste. Mocna własność Markowa mówi, że warunkowo
względem całej historii do $\tau$ dalszy proces jest
nową kopią procesu zaczynającego w stanie $0$.
Nie wolno zastąpić $\tau$ dowolną losową chwilą zależną
od przyszłych awarii: konieczny jest warunek czasu zatrzymania.
\end{przyklad}


\begin{zadanie}[Pierwszy rzeczywisty skok w trzech stanach]
Niech $E=\{0,1,2\}$, $q_{01}=2$, $q_{10}=1$,
$q_{12}=3$, a stan $2$ będzie pochłaniający.
Proces startuje w stanie $1$. Niech $J$ będzie chwilą
pierwszego skoku i $Y=X_J$. Wyznaczyć prawo $J$, prawo
$Y$ i obliczyć $\Prob(J\le1,Y=2)$.

\emph{Rozwiązanie.} Całkowita intensywność wyjścia ze
stanu $1$ wynosi $q_1=q_{10}+q_{12}=4$, wobec czego
$J\sim\operatorname{Exp}(4)$. Kierunek skoku jest
niezależny od $J$ i ma rozkład
$\Prob(Y=0)=1/4$, $\Prob(Y=2)=3/4$.
Miara $B\mapsto\Prob(J\in B,Y=2)$ na $[0,\infty)$
ma względem miary Lebesgue'a gęstość
\[
g(t)=q_{12}e^{-q_1t}=3e^{-4t},\qquad t\ge0.
\]
Zatem
$\Prob(J\le1,Y=2)=\int_0^1 3e^{-4t}\,\mathrm d\lambda_1(t)
=\tfrac34(1-e^{-4})$.
Nie należy tego mylić z $\Prob_1(X_1=2)$:
ten ostatni wynik obejmuje również ścieżki, na których
pierwszy skok prowadzi do $0$, a proces dopiero później
przechodzi przez $1$ do $2$.
\end{zadanie}

\section{Dwa stany: pełne rozwiązanie równań}
\begin{twierdzenie}[Przejścia modelu awarii i napraw]
Niech $E=\{0,1\}$, $\alpha,\beta>0$ oraz
$Q=\left(\begin{smallmatrix}-\alpha&\alpha\\
\beta&-\beta\end{smallmatrix}\right)$.
Wtedy dla $t\ge0$
\[
P(t)=\frac1{\alpha+\beta}
\begin{pmatrix}
\beta+\alpha e^{-(\alpha+\beta)t} &
\alpha(1-e^{-(\alpha+\beta)t})\\
\beta(1-e^{-(\alpha+\beta)t})&
\alpha+\beta e^{-(\alpha+\beta)t}
\end{pmatrix}.
\]
\end{twierdzenie}
\begin{proof}
Oznaczmy $u(t):=p_{01}(t)$, czyli prawdopodobieństwo
uszkodzenia w chwili $t$ przy starcie w stanie sprawnym.
Równanie w przód daje
$u'(t)=\alpha p_{00}(t)-\beta p_{01}(t)$.
Ponieważ suma elementów w pierwszym wierszu $P(t)$ jest $1$,
$p_{00}(t)=1-u(t)$, skąd
$u'(t)=\alpha-(\alpha+\beta)u(t)$ i $u(0)=0$.
Mnożąc obie strony równania liniowego przez czynnik
$e^{(\alpha+\beta)t}$, uzyskujemy
$(e^{(\alpha+\beta)t}u(t))'=\alpha e^{(\alpha+\beta)t}$.
Całkujemy od $0$ do $t$ i dostajemy
$u(t)=\alpha(1-e^{-(\alpha+\beta)t})/(\alpha+\beta)$.
Analogicznie, dla $v(t)=p_{10}(t)$ otrzymujemy
$v'(t)=\beta-(\alpha+\beta)v(t)$ i $v(0)=0$,
a więc $v(t)=\beta(1-e^{-(\alpha+\beta)t})/(\alpha+\beta)$.
Pozostałe wpisy wynikają z sum wierszy równych $1$.
\end{proof}

\begin{zadanie}[Prawdopodobieństwo awarii i informacja o przeszłości]
Urządzenie ma intensywność awarii $\alpha=2$ i naprawy
$\beta=1$ (w jednostkach: godzina$^{-1}$).
W chwili $0$ działa. Obliczyć prawdopodobieństwo,
że będzie uszkodzone po godzinie, a następnie
prawdopodobieństwo, że będzie uszkodzone po dwóch godzinach,
\emph{jeśli wiadomo, że działało po pierwszej godzinie}.

\emph{Rozwiązanie.} Z powyższego twierdzenia
$\Prob(X_1=1\mid X_0=0)=\tfrac23(1-e^{-3})$.
W drugim pytaniu warunkujemy względem stanu $X_1=0$,
a nie stanu początkowego. Własność Markowa daje
\[
\Prob(X_2=1\mid X_1=0,X_0=0)
=\Prob(X_1=1\mid X_0=0)=\frac23(1-e^{-3}).
\]
Gdybyśmy zamiast tego pytali bez warunku $X_1=0$
o stan po dwóch godzinach, odpowiedzią byłoby
$\tfrac23(1-e^{-6})$. To dwa różne pytania probabilistyczne.
\end{zadanie}

\section{Generatory, martyngały i wzór Dynkina}
\begin{definicja}[Generator skończonego łańcucha]
Dla skończonego $E$, macierzy intensywności $Q$ i funkcji
$f:E\to\R$ definiujemy
\[
(Lf)(i):=\sum_{j\in E}q_{ij}f(j)
=\sum_{j\ne i}q_{ij}\bigl(f(j)-f(i)\bigr),\qquad i\in E.
\]
Operator $L$ na przestrzeni funkcji na $E$ nazywamy
generatorem infinitezymalnym procesu. W zapisie wektorowym
$Lf=Qf$. Nazwa jest uzasadniona granicą
\[
Lf=\lim_{h\downarrow0}\frac{P(h)f-f}{h},
\]
w normie supremum: z szeregu wykładnika macierzowego
$P(h)=I+hQ+O(h^2)$ w normie operatorowej, więc
po zastosowaniu do $f$ i podzieleniu przez $h$
otrzymujemy wskazaną granicę.
\end{definicja}

Druga postać wyjaśnia sens operatora: z każdego stanu $i$
ważymy zmiany wartości $f(j)-f(i)$ intensywnościami
bezpośrednich skoków. Dla funkcji stałej generator znika,
co jest analogią do pochodnej funkcji stałej.

\begin{twierdzenie}[Wzór Dynkina i martyngał generatorowy]
\label{thm:dynkin}
Niech $E$ będzie skończone, $X$ łańcuchem w czasie
ciągłym o intensywnościach $Q$ i filtracji naturalnej.
Dla każdej funkcji $f:E\to\R$ proces
\[
M_t^f:=f(X_t)-f(X_0)-\int_0^t (Lf)(X_s)\,\mathrm d\lambda_1(s)
\]
jest całkowalnym martyngałem.
W szczególności dla każdego $t\ge0$
\[
\E f(X_t)=\E f(X_0)+
\int_0^t\E[(Lf)(X_s)]\,\mathrm d\lambda_1(s).
\]
Jeśli $\tau$ jest czasem zatrzymania względem tej filtracji,
ograniczonym punktowo przez pewną stałą,
to również
\[
\E f(X_\tau)=\E f(X_0)+
\E\int_0^\tau (Lf)(X_s)\,\mathrm d\lambda_1(s).
\]
\end{twierdzenie}
\begin{proof}
Wszystkie funkcje na skończonym $E$ są ograniczone.
Z równania w przód, w zapisie operatorowym,
$\frac{\dd}{\dd t}P(t)f=P(t)Lf$.
Całkując od $0$ do $u$ dostajemy
$P(u)f-f=\int_0^u P(v)Lf\,\mathrm d\lambda_1(v)$.
Dla $0\le s\le t$ własność Markowa implikuje
\[
\E[f(X_t)-f(X_s)\mid\calF_s]
=(P(t-s)f-f)(X_s).
\]
Jednocześnie $Lf$ jest funkcją ograniczoną na skończonym $E$, więc
$\int_s^t\E|Lf(X_u)|\,\mathrm d\lambda_1(u)\le(t-s)\norm{Lf}_\infty<\infty$.
Dla każdego $A\in\calF_s$ możemy więc zastosować twierdzenie
Fubiniego do $\E[\chi_A\int_s^t Lf(X_u)\,\mathrm d\lambda_1(u)]$,
a następnie własność Markowa pod całką. Charakterystyka
warunkowej wartości oczekiwanej daje
\[
\E\left[\int_s^t(Lf)(X_u)\,\mathrm d\lambda_1(u)\mid\calF_s\right]
=\int_0^{t-s}(P(v)Lf)(X_s)\,\mathrm d\lambda_1(v).
\]
Obie prawe strony są równe. Po odjęciu otrzymujemy
$\E[M_t^f-M_s^f\mid\calF_s]=0$.
Proces $M^f$ jest adaptowany: $X$ jest progresywnie
mierzalny, więc całka do chwili $t$ z ograniczonej
funkcji $(Lf)(X_s)$ jest $\calF_t^X$-mierzalna
na mocy twierdzenia Fubiniego o mierzalności całki.
Ponadto całka jest ciągła względem górnej granicy,
więc $M^f$ ma prawostronnie ciągłe trajektorie.
Zatem $M^f$ jest martyngałem. Całkowalność jest
natychmiastowa, ponieważ $f$ i $Lf$ są ograniczone,
a rozważamy skończony czas $t$.

Biorąc wartości oczekiwane uzyskujemy pierwszą tożsamość.
Dla $\tau\le T$ twierdzenie o próbkowaniu opcjonalnym
(twierdzenie~\ref{thm:optional-cont}) stosuje się do
prawostronnie ciągłego martyngału $M^f$, ponieważ
$\sup_{0\le t\le T}|M_t^f|\le2\norm{f}_\infty+T\norm{Lf}_\infty$.
Stąd $\E M_\tau^f=\E M_0^f=0$, co daje ostatni wzór.
\end{proof}

\begin{przyklad}[Przewidywana zmiana funkcji]
Dla modelu dwustanowego i $f(0)=0$, $f(1)=1$ mamy
$Lf(0)=\alpha$ i $Lf(1)=-\beta$.
Proces
\[
\chi_{\{X_t=1\}}-\chi_{\{X_0=1\}}
-\int_0^t\bigl(\alpha\chi_{\{X_s=0\}}
-\beta\chi_{\{X_s=1\}}\bigr)\,\mathrm d\lambda_1(s)
\]
jest martyngałem. Człon podcałkowy opisuje
\emph{chwilową przewidywaną zmianę} wskaźnika awarii,
a jego odjęcie usuwa dryf.
\end{przyklad}

\section{Rozkłady stacjonarne, równowaga szczegółowa i ergodyczność}
\begin{definicja}[Rozkład stacjonarny]
Dla nieeksplodującego łańcucha na skończonym lub
przeliczalnym $E$ z macierzą przejścia $P(t)$ rozkład probabilistyczny $\pi=(\pi_i)_{i\in E}$
nazywamy \emph{stacjonarnym}, jeżeli
$\pi P(t)=\pi$ dla każdego $t\ge0$.
Rozkład $\pi$ spełnia \emph{równania bilansu}, jeśli
$\pi Q=0$, przy czym na przeliczalnym $E$ wymagamy
istnienia odpowiednich sum w każdej współrzędnej.
\end{definicja}

\begin{twierdzenie}[Bilans i równowaga szczegółowa]
Niech $E$ będzie skończone i $Q$ macierzą intensywności.
Rozkład probabilistyczny $\pi$ jest stacjonarny wtedy i tylko
wtedy, gdy $\pi Q=0$. Warunkiem wystarczającym jest
\emph{równowaga szczegółowa}
\[
\pi_iq_{ij}=\pi_jq_{ji}\qquad\text{dla wszystkich }i,j\in E.
\]
\end{twierdzenie}
\begin{proof}
Jeśli $\pi P(t)=\pi$ dla każdego $t$, to przez różniczkowanie
w zerze otrzymujemy $\pi Q=0$.
Odwrotnie, gdy $\pi Q=0$, z rozwinięcia szeregu
$e^{tQ}=I+tQ+t^2Q^2/2!+\cdots$ wynika
$\pi e^{tQ}=\pi$, gdyż $\pi Q^k=0$ dla $k\ge1$.
Jeżeli zachodzi równowaga szczegółowa, to dla każdego $j$
\[
(\pi Q)_j=\sum_{i\ne j}\pi_iq_{ij}-\pi_j\sum_{k\ne j}q_{jk}
=\sum_{i\ne j}\pi_jq_{ji}-\pi_j\sum_{k\ne j}q_{jk}=0.
\]
Stąd $\pi Q=0$ i stacjonarność.
\end{proof}

\begin{stwierdzenie}[Równowaga szczegółowa a odwracalność]
Dla skończonego łańcucha o rozkładzie stacjonarnym
$\pi$ równowaga szczegółowa jest równoważna warunkowi
$\pi_i p_{ij}(t)=\pi_j p_{ji}(t)$ dla wszystkich $i,j$
i $t\ge0$. Przy starcie z rozkładu $\pi$ jest to
równoważne odwracalności w czasie: dla każdego $T>0$ i
$0\le t_1<\cdots<t_m\le T$ wektory
$(X_{t_1},\ldots,X_{t_m})$ oraz
$(X_{T-t_1},\ldots,X_{T-t_m})$ mają ten sam rozkład.
\end{stwierdzenie}
\begin{proof}
Niech $D=\operatorname{diag}(\pi_i:i\in E)$.
Równowaga szczegółowa oznacza $DQ=Q^{\mathsf T}D$.
Przez indukcję $DQ^n=(Q^{\mathsf T})^nD$, zatem
z szeregu wykładnika macierzowego
$DP(t)=P(t)^{\mathsf T}D$. Odwrotnie, różniczkowanie
tej tożsamości w zerze daje $DQ=Q^{\mathsf T}D$.
Nie trzeba zakładać odwracalności macierzy $D$.
Dla ustalonego ciągu stanów $i_1,\ldots,i_m$,
ze stacjonarności i własności Markowa mamy
\[
\pi_{i_1}\prod_{k=1}^{m-1}p_{i_k i_{k+1}}(t_{k+1}-t_k)
=\pi_{i_m}\prod_{k=1}^{m-1}p_{i_{k+1} i_k}(t_{k+1}-t_k).
\]
Równość otrzymujemy, kolejno stosując bilans szczegółowy
dla $P(t_{k+1}-t_k)$, bez dzielenia przez $\pi_i$.
Prawa strona jest prawdopodobieństwem odpowiedniego
zdarzenia dla wektora odwróconych obserwacji.
Odwrotnie, odwracalność dla pary chwil $0,t$ przy
$T=t>0$ daje $\pi_i p_{ij}(t)=\pi_j p_{ji}(t)$.
\end{proof}

Nie należy mylić bilansu całkowitego z równowagą szczegółową:
warunek szczegółowy jest silniejszy i nie zachodzi na przykład
dla jednokierunkowego ruchu po cyklu z co najmniej trzema
stanami, mimo istnienia rozkładu stacjonarnego.

\begin{twierdzenie}[Granica w modelu skończonym]
Niech $E$ będzie skończonym zbiorem stanów, $Q$ macierzą
intensywności, a graf przejść z krawędzią $i\to j$ wtedy
i tylko wtedy, gdy $q_{ij}>0$, niech będzie silnie spójny.
Wtedy istnieje dokładnie jeden rozkład stacjonarny $\pi$;
wszystkie jego współrzędne są dodatnie i dla każdego
$i,j\in E$ zachodzi
$\lim_{t\to\infty}p_{ij}(t)=\pi_j$.
\end{twierdzenie}
\begin{proof}
Weźmy $\Lambda>\max_i q_i$ i $R=I+Q/\Lambda$.
Z silnej spójności grafu intensywności wynika
nieprzywiedlność dyskretnego łańcucha o macierzy $R$,
gdyż $R_{ij}>0$ dla $i\ne j$ dokładnie na dozwolonych
krawędziach. Ponieważ $R_{ii}>0$ dla każdego $i$,
łańcuch ten jest nieokresowy.
Twierdzenie o zbieżności skończonych nieprzywiedlnych
i nieokresowych łańcuchów z tomu I zapewnia istnienie
jedynego stacjonarnego $\pi$, dodatniego w każdym stanie,
i $R^n_{ij}\to\pi_j$ dla wszystkich $i,j$.
Warunek $\pi R=\pi$ jest równoważny $\pi Q=0$,
a więc rozkład jest też stacjonarny dla czasu ciągłego.

Z uniformizacji mamy
$p_{ij}(t)=\sum_{n\ge0}\Prob(N_t=n)R^n_{ij}$,
gdzie $N_t\sim\operatorname{Poisson}(\Lambda t)$.
Dla $\varepsilon>0$ wybierzmy $n_0$ tak, by
$|R^n_{ij}-\pi_j|<\varepsilon$ dla $n\ge n_0$.
Wtedy
\[
|p_{ij}(t)-\pi_j|
\le\Prob(N_t<n_0)+\varepsilon.
\]
Ponieważ $N_t\to\infty$ według prawdopodobieństwa
(gdy $t\to\infty$, choćby z nierówności Czebyszewa),
pierwszy składnik dąży do zera. Dajemy $\varepsilon\downarrow0$.
Jedyność rozkładu stacjonarnego wynika także z tego,
że każdy rozkład stacjonarny dla $Q$ jest stacjonarny dla $R$.
\end{proof}

\begin{twierdzenie}[Ergodyczne średnie czasowe w modelu skończonym]
Niech $X$ będzie łańcuchem na skończonej przestrzeni
stanów $E$, o nieprzywiedlnej macierzy intensywności
i jedynym rozkładzie stacjonarnym $\pi$.
Dla każdej funkcji
$f:E\to\R$ i dowolnego rozkładu początkowego
\[
\lim_{T\to\infty}\frac1T\int_0^T f(X_s)\,
\mathrm d\lambda_1(s)=\sum_{i\in E}\pi_i f(i)
\quad\text{prawie na pewno i w }L^2.
\]
\end{twierdzenie}
\begin{proof}
Najpierw uzasadnimy jednostajne tempo zbieżności
macierzy przejścia. Nieprzywiedlność i uniformizacja
dają $p_{ij}(1)>0$ dla wszystkich $i,j$: istnieje
$n$ z $(R^n)_{ij}>0$, a jego waga Poissona jest dodatnia.
Niech $d=|E|$ i
$\varepsilon=\tfrac12 d\min_{i,j}p_{ij}(1)\in(0,1)$.
Możemy zapisać $P(1)=\varepsilon U+(1-\varepsilon)S$,
gdzie każdy wiersz $U$ jest rozkładem jednostajnym
na $E$, a $S$ jest macierzą stochastyczną.
Dla rozkładów wierszowych $\nu,\eta$ mamy
$(\nu-\eta)U=0$ oraz
$\norm{(\nu-\eta)S}_1\le\norm{\nu-\eta}_1$
z nierówności trójkąta i sum wierszy $S$ równych $1$.
Stąd, używając półgrupy przejścia i stacjonarności,
\[
\norm{\delta_i P(t)-\pi}_1\le
2(1-\varepsilon)^{\lfloor t\rfloor},\qquad t\ge0,
\]
gdzie $\delta_i$ jest rozkładem skupionym w $i$;
dla ułamkowej części czasu stosujemy kontraktywność
macierzy stochastycznej w normie $\ell^1$.

Niech $g=f-\sum_i\pi_i f(i)$. Z powyższego wynika
$\norm{P(t)g}_\infty\le
2\norm{g}_\infty(1-\varepsilon)^{\lfloor t\rfloor}$.
Dla $0\le u\le v$ własność Markowa daje
$\E[g(X_u)g(X_v)]=\E[g(X_u)(P(v-u)g)(X_u)]$.
Wobec ograniczoności $g$ i wspólnej mierzalności
procesu (wynikającej z progresywnej mierzalności)
możemy zastosować twierdzenie Fubiniego.
Całka po kwadracie $[0,T]^2$ jest dwukrotnością
całki po trójkącie $0\le u<v\le T$; przekątna ma
miarę Lebesgue'a zero. Po podstawieniu $w=v-u$
otrzymujemy
\[
\E\left|\frac1T\int_0^T g(X_s)\,\mathrm d\lambda_1(s)\right|^2
\le \frac{4\norm{g}_\infty^2}{T}
\int_0^\infty(1-\varepsilon)^{\lfloor w\rfloor}
\,\mathrm d\lambda_1(w)
=\frac{4\norm{g}_\infty^2}{\varepsilon T}.
\]
Otrzymujemy zbieżność w $L^2$.
Dla $T_n=n^2$ nierówność Czebyszewa i pierwszy
lemat Borela--Cantellego dają zbieżność prawie na pewno
wzdłuż tego ciągu (stosujemy je dla progów $1/k$,
$k\ge1$). Dla $n^2\le T<(n+1)^2$ różnica średnich
w chwilach $T$ i $n^2$ ma moduł nie większy niż
$2\norm{g}_\infty(2n+1)/n^2$, co dąży do zera.
To kończy dowód dla wszystkich $T\to\infty$.
\end{proof}

\begin{przyklad}[Łańcuch trzystanowy na cyklu]
Niech $E=\{0,1,2\}$ i jedyne przejścia mają postać
$0\to1$, $1\to2$, $2\to0$, każde z intensywnością
$\lambda>0$. Wtedy
\[
Q=\lambda\begin{pmatrix}-1&1&0\\0&-1&1\\1&0&-1\end{pmatrix}.
\]
Rozkład $\pi=(1/3,1/3,1/3)$ spełnia $\pi Q=0$,
więc jest stacjonarny. Graf jest silnie spójny,
zatem proces zbiega rozkładem do $\pi$.
Jednak równowaga szczegółowa nie zachodzi, bo
$\pi_0q_{01}=\lambda/3>0$ i $\pi_1q_{10}=0$.
Proces ma więc stacjonarny, niezerowy przepływ wokół cyklu.
\end{przyklad}

\section{Łańcuchy urodzin i śmierci}
\begin{definicja}[Proces urodzin i śmierci]
Niech $E=\N_0$. Procesem urodzin i śmierci nazywamy
minimalny łańcuch Markowa o jedynych dozwolonych przejściach
$n\to n+1$ z intensywnością $\lambda_n\ge0$ i
$n\to n-1$ z intensywnością $\mu_n\ge0$ dla $n\ge1$,
przy $\mu_0:=0$. Pozostałe wyrazy pozadiagonalne są zerowe,
a $q_{nn}=-(\lambda_n+\mu_n)$.
\end{definicja}

Nazwa nawiązuje do populacji, w której każda zmiana liczebności
jest zwiększeniem lub zmniejszeniem o jeden. Ta sama
matematyka opisuje jednak liczbę klientów w kolejce, cząstek
w systemie czy liczebność aktywnych urządzeń.

\begin{stwierdzenie}[Rozkład stacjonarny przy równowadze lokalnej]
Załóżmy, że $\lambda_n>0$ dla $n\ge0$,
$\mu_n>0$ dla $n\ge1$, oraz że proces jest nieeksplodujący.
Jeżeli
\[
Z:=1+\sum_{n=1}^\infty
\prod_{k=0}^{n-1}\frac{\lambda_k}{\mu_{k+1}}<\infty,
\]
to liczby
$\pi_0:=1/Z$ i
$\pi_n:=Z^{-1}\prod_{k=0}^{n-1}\lambda_k/\mu_{k+1}$
dla $n\ge1$ tworzą rozkład probabilistyczny spełniający
równowagę szczegółową
$\pi_n\lambda_n=\pi_{n+1}\mu_{n+1}$.
W szczególności $\pi$ jest rozkładem stacjonarnym procesu.
Odwrotnie, istnienie rozkładu stacjonarnego wymaga $Z<\infty$,
a powyższy rozkład jest wówczas jedyny.
\end{stwierdzenie}
\begin{proof}
Wszystkie $\pi_n$ są nieujemne, a ich suma z definicji
jest równa $Z/Z=1$. Przez skrócenie kolejnych iloczynów
otrzymujemy $\pi_n\lambda_n=\pi_{n+1}\mu_{n+1}$.
Wykazaliśmy zatem równowagę szczegółową.
Pozostaje udowodnić stacjonarność \emph{dla procesu
o przeliczalnej przestrzeni stanów}; nie wystarczy tylko zapisać
formalnej równości $\pi Q=0$.

Dla $K\in\N$ rozważmy łańcuch na $\{0,1,\ldots,K\}$
z tymi samymi intensywnościami, z wyjątkiem tego,
że tłumimy przejście $K\to K+1$.
Jest to skończony, nieeksplodujący proces,
a jego rozkład
\[
\pi_n^{(K)}:=\frac{\pi_n}{\sum_{k=0}^K\pi_k},
\qquad 0\le n\le K,
\]
spełnia równowagę szczegółową dla wszystkich jego krawędzi,
a więc jest stacjonarny w myśl twierdzenia skończonego.

Dla ustalonego stanu początkowego $i\in\N_0$ skonstruujmy
sprzężenie procesu pierwotnego i obciętego przy użyciu
wspólnych zegarów potencjalnych przejść w górę i w dół.
W stanie $n$ zegary mają niezależne czasy wykładnicze
o parametrach $\lambda_n,\mu_n$ (parametr zero oznacza
czas nieskończony). Minimum ma parametr sumy, a kierunek
jest niezależny od minimum, z prawdopodobieństwami
proporcjonalnymi do parametrów, co wynika z iloczynu
funkcji przeżycia i gęstości. W stanie $K$ proces
obcięty ignoruje zegar przejścia w górę; zachowuje
zegar przejścia w dół. Do pierwszego wyjścia procesu pierwotnego poza
$\{0,\ldots,K\}$ oba procesy są jednakowe (bez
samej chwili wyjścia). Po rozdzieleniu proces
obcięty może dokończyć oczekiwanie według już
wylosowanego zegara w dół; jego czas przebywania
od wejścia do $K$ jest więc dokładnie wykładniczy
z parametrem $\mu_K$. Po zakończeniu tego pobytu
używamy nowych niezależnych zegarów dla procesu
obciętego. W ten sposób dla każdego $K$ mamy sprzężenie
o właściwych rozkładach brzegowych. Ponieważ oryginalny
proces jest nieeksplodujący,
na każdym $[0,t]$ wykonuje tylko skończenie wiele skoków
prawie na pewno; jego trajektoria ma zatem na $[0,t]$
skończone maksimum. W opisanym sprzężeniu różnica
prawdopodobieństw przejścia jest oszacowana przez
$\Prob_i(\sup_{0\le v\le t}X_v>K)$, które dąży do zera.
Stąd dla każdego $i,j$ i $t$
\[
P_t^{(K)}(i,j)\longrightarrow P_t(i,j)
\quad (K\to\infty,\ K\ge i),
\]
gdzie $P_t^{(K)}$ oznacza macierz przejścia łańcucha
obciętego.

Teraz ustalmy $j$ i $t$. Stacjonarność $\pi^{(K)}$ daje
\[
\pi_j^{(K)}=\sum_{i=0}^K\pi_i^{(K)}P_t^{(K)}(i,j).
\]
Dla ustalonego $i$ mamy $\pi_i^{(K)}\to\pi_i$
i zbieżność odpowiednich prawdopodobieństw przejścia.
Ponieważ $\sum_{i\ge0}\pi_i=1$, ogon
$\sum_{i>M}\pi_i$ jest dowolnie mały, a dla dużych $K$
$\pi_i^{(K)}\le2\pi_i$. Stąd możemy najpierw przepuścić
$K\to\infty$ w skończonej sumie po $i\le M$,
a następnie $M\to\infty$ kontrolując ogon przez
$2\sum_{i>M}\pi_i$. Otrzymujemy
$\pi_j=\sum_{i\ge0}\pi_iP_t(i,j)$.
Ponieważ $j,t$ były dowolne, $\pi$ jest stacjonarna.

Dla konieczności niech $\nu$ będzie dowolnym rozkładem
stacjonarnym. Ustalmy $n\ge0$ i funkcję
$f=\chi_{\{0,\ldots,n\}}$ na $\N_0$.
Jej generator, liczony za pomocą macierzy intensywności
procesu urodzin i śmierci, jest ograniczony:
\[
g(i)=
\begin{cases}
-\lambda_n,&i=n,\\
\mu_{n+1},&i=n+1,\\
0,&\text{w pozostałych stanach}.
\end{cases}
\]
Dla każdego $K\ge n+1$ generator skończonego procesu
obciętego, zastosowany do ograniczenia $f$, jest
ograniczeniem tej samej funkcji $g$. Ustalmy stan
początkowy $i$ i bierzmy $K\ge\max(i,n+1)$.
Wzór Dynkina (twierdzenie~\ref{thm:dynkin}) daje
\[
\E_i f(X_t^{(K)})-f(i)
=\int_0^t\E_i g(X_s^{(K)})\,\mathrm d\lambda_1(s).
\]
Z powyższego sprzężenia wynika zbieżność wartości
oczekiwanych każdej ograniczonej funkcji przy ustalonej
chwili $s$: moduł różnicy jest ograniczony przez
dwukrotność jej normy supremum pomnożoną przez
$\Prob_i(\sup_{0\le v\le s}X_v>K)$.
Stosując twierdzenie o zbieżności zmajoryzowanej
pod całką, otrzymujemy wzór Dynkina dla tej konkretnej
funkcji $f$ i procesu nieobciętego.
Uśredniamy teraz względem $\nu$; zamiana sumy z całką
jest uzasadniona ograniczonością $g$ i $\sum_i\nu_i=1$.
Stacjonarność daje dla $t>0$
\[
0=t(-\nu_n\lambda_n+\nu_{n+1}\mu_{n+1}).
\]
Zatem $\nu$ spełnia równowagę szczegółową na każdej
krawędzi, a indukcja prowadzi do
$\nu_n=\nu_0\prod_{k=0}^{n-1}\lambda_k/\mu_{k+1}$.
Gdyby $\nu_0=0$, wszystkie współrzędne byłyby zerowe.
Normalizacja wymaga więc $Z<\infty$ oraz $\nu_0=1/Z$,
co dowodzi konieczności i jedyności. Nie zakładaliśmy
ograniczoności wszystkich intensywności.
\end{proof}

\begin{uwaga}[Wersja nieskończona wymaga ostrożności]
W poprzednim stwierdzeniu założyliśmy nieeksplodowanie.
W ogólnych modelach na przestrzeni przeliczalnej
formalna równość $\pi Q=0$ nie jest automatycznie
równoważna stacjonarności minimalnego procesu,
gdy masa może uciekać przez eksplozję. Z tego powodu
nie przenosimy bez komentarza wyników z przypadku
skończonego na wszystkie macierze $Q$.
\end{uwaga}


\begin{uwaga}[Dlaczego obcięcie nieskończonego łańcucha wymaga sprzężenia?]
W dowodzie stacjonarności procesu urodzin i śmierci
przybliżamy łańcuch na $\N_0$ przez łańcuchy na
$\{0,\ldots,K\}$. Nie wystarczy tylko stwierdzić, że
macierze intensywności są podobne: w stanie brzegowym $K$
mają inne czasy przebywania. Można jednak \emph{sprzęgnąć}
procesy, korzystając z tych samych niezależnych zegarów
potencjalnych przejść $n\to n+1$ i $n\to n-1$.
W łańcuchu obciętym ignorujemy wyłącznie impulsy
$K\to K+1$; pozostałe są wspólne dla obu procesów.
Do pierwszego wejścia oryginalnego procesu w $K+1$
trajektorie obu procesów są więc identyczne.
Przy założonym braku eksplozji liczba skoków procesu
oryginalnego na $[0,t]$ jest prawie na pewno skończona,
co oznacza, że na tym przedziale osiąga tylko skończenie
wiele stanów. Prawdopodobieństwo rozbieżności trajektorii w tym
sprzężeniu na $[0,t]$ nie przekracza
$\Prob_i(\sup_{0\le v\le t}X_v>K)$ i dąży do zera.
Nie potrzebujemy zakładać wspólnej realizacji wszystkich
przybliżeń dla różnych $K$.
To uzasadnia zbieżność ich praw przejścia, a dopiero
następnie można przejść do granicy w równaniu stacjonarności.
\end{uwaga}

\begin{uwaga}[Zanik pamięci a własność Markowa]
Własność Markowa mówi, że po znajomości $X_t$ cała
wcześniejsza historia nie jest potrzebna do prognozowania
przyszłości. Nie oznacza, że rozkład $X_t$ przestaje
zależeć od stanu początkowego w krótkim czasie.
Asymptotyczny zanik tej zależności wynika w modelu
skończonym z nieprzywiedlności i twierdzenia o granicy
$P(t)$; dla niektórych procesów na przestrzeniach
nieskończonych taki wynik może być fałszywy.
\end{uwaga}

\section{Model kolejki \texorpdfstring{$M/M/1$}{M/M/1}}
\begin{przyklad}[Pojedynczy serwer]
Klienci napływają zgodnie z procesem Poissona o intensywności
$\lambda>0$. Czasy obsługi są niezależne, mają rozkład
$\operatorname{Exp}(\mu)$ z $\mu>0$ i są niezależne
od napływu. Jeden serwer obsługuje co najwyżej jednego klienta
naraz, a klienci czekają bez ograniczenia pojemności.
Stanem $X_t$ jest liczba osób w systemie (obsługiwanych
i oczekujących). Przy $X_t=n$ nowy klient zwiększa stan do
$n+1$ z intensywnością $\lambda$, a przy $n\ge1$
zakończenie obsługi zmniejsza stan do $n-1$ z intensywnością
$\mu$. Jest to proces urodzin i śmierci z
$\lambda_n=\lambda$ i $\mu_n=\mu$ dla $n\ge1$.

Na mocy twierdzenia~\ref{thm:bounded-rates} proces
jest nieeksplodujący, ponieważ $q_n\le\lambda+\mu$.
Z poprzedniego stwierdzenia otrzymujemy dla
$\rho:=\lambda/\mu<1$:
\[
\pi_n=(1-\rho)\rho^n,\qquad n\in\N_0.
\]
Konieczność warunku $\rho<1$ oraz jedyność rozkładu
wynikają już z ogólnego kryterium: tutaj $Z=\sum_{n\ge0}\rho^n$.
Dla ilustracji można też wyprowadzić równania bilansu
bezpośrednio, korzystając z ograniczoności intensywności.
Każdy rozkład stacjonarny spełnia $\pi Q=0$: dla
$h\downarrow0$ prawdopodobieństwa przejścia spełniają
$p_{ij}(h)=\chi_{\{i=j\}}+hq_{ij}+O(h^2)$,
z błędem jednostajnym względem $i,j$.
Istotnie, uniformizacja z zegarem o intensywności
$\Lambda>\lambda+\mu$ działa także w tym modelu
przeliczalnym: $R=I+Q/\Lambda$ jest stochastyczna,
a prawdopodobieństwo co najmniej dwóch impulsów jest
$O(h^2)$, niezależnie od stanu początkowego.
Rozkład stacjonarny spełnia $\pi P(h)=\pi$,
więc po odjęciu $\pi$, podzieleniu przez $h$ i przejściu
do granicy można sumować po $i$ dzięki jednostajnemu
oszacowaniu błędu i $\sum_i\pi_i=1$.
Równanie bilansu w stanie $0$ brzmi:
$\lambda\pi_0=\mu\pi_1$.
Bilans w stanie $n\ge1$ brzmi
$\lambda\pi_{n-1}+\mu\pi_{n+1}=(\lambda+\mu)\pi_n$.
Indukcyjnie wynika stąd $\pi_{n+1}=\rho\pi_n$, gdyż
równość dla $n=0$ jest początkiem indukcji, a podstawienie
$\lambda\pi_{n-1}=\mu\pi_n$ do równania dla $n$
prowadzi do $\mu\pi_{n+1}=\lambda\pi_n$.
Stąd $\pi_n=\rho^n\pi_0$. Musi być $\pi_0>0$,
gdyż w przeciwnym razie wszystkie $\pi_n$ byłyby zerowe.
Możliwość normalizacji wymaga zatem $\rho<1$. Wtedy $\pi_0=1-\rho$.
W tym stanie równowagi
$\Prob(X_t=0)=1-\rho$ i
$\E X_t=\sum_{n\ge0}n(1-\rho)\rho^n=\rho/(1-\rho)$.
Ostatni wzór otrzymujemy różniczkując szereg geometryczny
$\sum_{n\ge0}\rho^n=(1-\rho)^{-1}$ dla $0<\rho<1$.

Przy $\rho\ge1$
proces nadal istnieje, lecz nie ma rozkładu stacjonarnego
normalizowalnego na $\N_0$.
\end{przyklad}

\begin{zadanie}[Kolejka i liczba klientów]
Do stanowiska przychodzi średnio $3$ klientów na godzinę,
a jeden serwer obsługuje średnio $5$ klientów na godzinę;
napływ i czasy obsługi spełniają założenia modelu $M/M/1$.
Obliczyć w stanie stacjonarnym: (a) prawdopodobieństwo
pustego systemu; (b) prawdopodobieństwo obecności co
najmniej trzech klientów; (c) średnią liczbę klientów.

\emph{Rozwiązanie.} Mamy $\rho=3/5$. Zatem
$\pi_0=1-\rho=2/5$.
Dalej
$\Prob(X\ge3)=\sum_{n=3}^\infty(1-\rho)\rho^n
=\rho^3=27/125$.
Średnia wynosi $\rho/(1-\rho)=\frac{3/5}{2/5}=3/2$.
Wszystkie te wyniki dotyczą \emph{stanu stacjonarnego};
nie są poprawnymi odpowiedziami na pytania o układ,
który właśnie został uruchomiony z zerową liczbą klientów.
\end{zadanie}

\section{Model Ehrenfesta w czasie ciągłym}
\begin{przyklad}[Dwie komory]
Niech $K\ge1$ będzie liczbą cząstek niezależnie
przemieszczających się
między komorami $A$ i $B$. Każda cząstka w $B$ przechodzi
do $A$ z intensywnością $\alpha>0$, a każda cząstka w $A$
przechodzi do $B$ z intensywnością $\beta>0$.
Niech $X_t$ oznacza liczbę cząstek w $A$ w chwili $t$.
Gdy $X_t=i$, w $B$ znajduje się $K-i$ cząstek, więc
$\lambda_i=(K-i)\alpha$, a $\mu_i=i\beta$.
Istotnie, minimum niezależnych wykładniczych czasów
do zmian komór poszczególnych cząstek ma intensywność
$(K-i)\alpha+i\beta$, a prawdopodobieństwa kierunków
są proporcjonalne do odpowiednich sum intensywności.
Po zmianie stanu pozostałe zegary zachowują
własność braku pamięci. Prawo następnego czasu
i kierunku zależy więc jedynie od liczby $i$,
nie od tożsamości cząstek w komorze.
Łańcuch na $\{0,\ldots,K\}$ jest nieprzywiedlny.

Dla $p:=\alpha/(\alpha+\beta)$ rozkład
\[
\pi_i=\binom Ki p^i(1-p)^{K-i},\qquad 0\le i\le K,
\]
jest stacjonarny. Rzeczywiście, dla $0\le i<K$
\begin{align*}
\frac{\pi_{i+1}}{\pi_i}
&=\frac{K-i}{i+1}\frac{p}{1-p}
=\frac{(K-i)\alpha}{(i+1)\beta}
=\frac{\lambda_i}{\mu_{i+1}},
\end{align*}
co dokładnie oznacza równowagę szczegółową
$\pi_i\lambda_i=\pi_{i+1}\mu_{i+1}$.
Rozkład dwumianowy pojawia się zatem nie tylko
przy niezależnych próbach, lecz także jako prawo
stacjonarne łańcucha w czasie ciągłym.
\end{przyklad}

\section{Czas pochłonięcia i macierz fundamentalna}
\begin{twierdzenie}[Oczekiwany czas pochłonięcia]
\label{thm:ctmc-absorption}
Niech $E=T\mathbin{\dot\cup}A$ będzie skończonym zbiorem stanów,
$A\ne\varnothing$ zbiorem stanów pochłaniających,
a $T\ne\varnothing$ zbiorem stanów niepochłaniających. Przy uporządkowaniu
najpierw stanów $T$ macierz intensywności ma blokową postać
\[
Q=\begin{pmatrix}B&C\\0&0\end{pmatrix}.
\]
Załóżmy, że startując z każdego $i\in T$, proces
osiąga $A$ z prawdopodobieństwem jeden.
Wtedy $B$ jest odwracalna, macierz
$G:=(-B)^{-1}$ ma nieujemne wyrazy i
\[
G_{ij}=\E_i\left[\int_0^\tau\chi_{\{X_t=j\}}\,\mathrm d\lambda_1(t)\right],
\qquad
\E_i\tau=\sum_{j\in T}G_{ij},
\]
gdzie $\tau=\inf\{t\ge0:X_t\in A\}$, a
$\E_i$ oznacza wartość oczekiwaną dla procesu
startującego ze stanu $i$.
\end{twierdzenie}
\begin{proof}
Stany z $A$ są pochłaniające, dlatego dla $i,j\in T$
$\Prob_i(X_t=j,\tau>t)=\Prob_i(X_t=j)$.
Macierz przejścia ograniczona do $T$ jest równa $e^{tB}$:
wynika to z blokowo-trójkątnej postaci $e^{tQ}$,
którą widać bezpośrednio z potęg macierzy $Q$.
Stąd całkując funkcję wskaźnikową i stosując Tonellego,
\[
\E_i\int_0^\tau\chi_{\{X_t=j\}}\,\mathrm d\lambda_1(t)
=\int_0^\infty\Prob_i(X_t=j,\tau>t)\,\mathrm d\lambda_1(t)
=\int_0^\infty(e^{tB})_{ij}\,\mathrm d\lambda_1(t).
\]

Musimy wykazać zbieżność tej całki oraz odwracalność $B$.
Ponieważ proces z każdego stanu $i\in T$ zostaje pochłonięty
prawie na pewno, istnieje $t_i>0$ takie, że
$\Prob_i(\tau>t_i)<1/2$. Biorąc $t_0=\max_i t_i$,
z monotoniczności funkcji przeżycia dostajemy
$\sup_{i\in T}\Prob_i(\tau>t_0)\le1/2$.
Ponieważ $A$ jest pochłaniający, $\{\tau>t\}$
jest równoważne $\{X_t\in T\}$. Własność Markowa
w chwilach $nt_0$ daje dla $i\in T$
\[
\Prob_i(\tau>(n+1)t_0)
=\E_i[\chi_{\{\tau>nt_0\}}\Prob_{X_{nt_0}}(\tau>t_0)]
\le\tfrac12\Prob_i(\tau>nt_0).
\]
Indukcyjnie $\sup_{i\in T}\Prob_i(\tau>nt_0)\le2^{-n}$.
Zatem wierszowe sumy macierzy $e^{tB}$ maleją
wykładniczo do zera, a nieujemne wyrazy tej macierzy
są całkowalne na $[0,\infty)$. Niech
$G=\int_0^\infty e^{tB}\,\mathrm d\lambda_1(t)$.
Z równania $\frac{\dd}{\dd t}e^{tB}=Be^{tB}$ wynika
\[
BG=\lim_{t\to\infty}e^{tB}-I=-I.
\]
W skończonym wymiarze istnienie prawej odwrotności
macierzy kwadratowej oznacza odwracalność, a
$G=(-B)^{-1}$. Nieujemność wyrazów jest widoczna
z całkowej reprezentacji. Korzystając z tożsamości
$\int_0^\tau\sum_{j\in T}\chi_{\{X_t=j\}}\,\mathrm d\lambda_1(t)=\tau$
(na $[0,\tau)$ proces pozostaje w $T$)
i liniowości wartości oczekiwanej dla skończonej sumy,
otrzymujemy drugi wzór.
\end{proof}

\begin{zadanie}[Pełne obliczenie czasu pochłonięcia]
Rozważmy łańcuch z trzema stanami i intensywnościami
$0\to1$ z intensywnością $2$, $1\to0$ z intensywnością $1$
oraz $1\to2$ z intensywnością $3$.
Stan $2$ jest pochłaniający. Obliczyć oczekiwany czas
osiągnięcia stanu $2$ przy starcie ze stanu $0$ i $1$,
a także oczekiwany łączny czas spędzony w stanach $0$ i $1$
przed pochłonięciem przy starcie z $0$.

\emph{Rozwiązanie.} Macierz intensywności jest równa
\[
Q=\begin{pmatrix}-2&2&0\\1&-4&3\\0&0&0\end{pmatrix},
\quad
B=\begin{pmatrix}-2&2\\1&-4\end{pmatrix}.
\]
Pochłonięcie zachodzi prawie na pewno: przy każdym wejściu
do stanu $1$ następny skok prowadzi do $2$ z
prawdopodobieństwem $3/4$, a powroty do stanu $1$ po skoku
do $0$ następują prawie na pewno. Prawdopodobieństwo
uniknięcia stanu $2$ po $m$ wizytach w $1$ wynosi $(1/4)^m$,
co dąży do zera.
Z twierdzenia~\ref{thm:ctmc-absorption}
\[
G=(-B)^{-1}
=\begin{pmatrix}2&-2\\-1&4\end{pmatrix}^{-1}
=\frac16\begin{pmatrix}4&2\\1&2\end{pmatrix}.
\]
Przy starcie z $0$ oczekiwane czasy zajmowania kolejno
stanu $0$ i stanu $1$ wynoszą $4/6=2/3$ oraz $2/6=1/3$.
Ich suma to $\E_0\tau=1$. Przy starcie z $1$ suma
elementów drugiego wiersza daje $\E_1\tau=1/2$.

\emph{Sprawdzenie bez macierzy.} Oznaczmy
$u_i=\E_i\tau$, $u_2=0$.
Ze stanu $0$ oczekujemy średnio $1/2$ jednostki czasu
i przechodzimy do $1$, więc $u_0=1/2+u_1$.
Ze stanu $1$ czekamy średnio $1/4$, po czym z
prawdopodobieństwem $1/4$ przechodzimy do $0$,
a z $3/4$ do $2$. Stąd
$u_1=1/4+(1/4)u_0$. Rozwiązanie daje
$u_0=1$, $u_1=1/2$ zgodnie z rachunkiem macierzowym.
To jest praktyczna metoda \emph{analizy pierwszego skoku}.
\end{zadanie}

\section{Jak czytać prawdopodobieństwa trajektorii?}
Przy procesie dyskretnym używaliśmy wzoru
$\Prob(X_0=i_0,\ldots,X_n=i_n)$.
W czasie ciągłym należy rozróżnić obserwacje
w \emph{ustalonych chwilach} od zarejestrowania
\emph{dokładnych czasów skoków}.

\begin{stwierdzenie}[Wspólny rozkład obserwacji w ustalonych chwilach]
Niech $X$ będzie skończonym łańcuchem w czasie ciągłym,
$0=t_0<t_1<\cdots<t_n$, a $\mu$ rozkładem początkowym.
Wtedy dla każdego $i_0,\ldots,i_n\in E$
\[
\Prob(X_{t_0}=i_0,\ldots,X_{t_n}=i_n)
=\mu_{i_0}\prod_{k=1}^n
 p_{i_{k-1},i_k}(t_k-t_{k-1}).
\]
\end{stwierdzenie}
\begin{proof}
Jeżeli prawdopodobieństwo wcześniejszej części zdarzenia
jest dodatnie, stosujemy łańcuchowy wzór na prawdopodobieństwo
warunkowe. Własność Markowa daje
\[
\Prob(X_{t_k}=i_k\mid X_{t_0}=i_0,\ldots,X_{t_{k-1}}=i_{k-1})
=p_{i_{k-1},i_k}(t_k-t_{k-1}).
\]
Mnożymy kolejne czynniki, rozpoczynając od
$\Prob(X_{t_0}=i_0)=\mu_{i_0}$.
Jeżeli prawdopodobieństwo wcześniejszego zdarzenia jest zero,
to prawdopodobieństwo całej ścieżki obserwacji również
jest zero, a iloczyn po prawej zawiera zerowy czynnik
na pierwszym etapie, na którym otrzymujemy zerowy iloczyn.
W ten sposób wzór obejmuje także zdarzenia zerowe.
\end{proof}

\begin{uwaga}[Prawdopodobieństwo dokładnej chwili skoku]
Jeżeli $J_1$ jest pierwszą chwilą skoku ze stanu $i$ o
$q_i>0$, to $J_1\sim\operatorname{Exp}(q_i)$ i
$\Prob(J_1=t)=0$ dla każdego ustalonego $t\ge0$.
Nie pytamy więc o \emph{prawdopodobieństwo} dokładnego czasu,
lecz o jego \emph{gęstość}. Dla ściśle rosnących
$0<s_1<\cdots<s_m<t$ oraz ciągu rzeczywistych zmian
$i_0\to i_1\to\cdots\to i_m$ gęstość zdarzenia:
``skoki nastąpiły w chwilach $s_1,\ldots,s_m$,
a do $t$ nie nastąpił następny skok'' jest równa
\[
\mu_{i_0}\left[\prod_{k=1}^m
q_{i_{k-1},i_k}e^{-q_{i_{k-1}}(s_k-s_{k-1})}\right]
e^{-q_{i_m}(t-s_m)},\qquad s_0:=0.
\]
Wynika to z mnożenia gęstości wykładniczych czasów
przebywania i prawdopodobieństw wyboru kolejnych stanów:
$[q_i e^{-q_i u}](q_{ij}/q_i)=q_{ij}e^{-q_i u}$.
To zapis, od którego zaczyna się matematyczne formułowanie
wiarygodności ciągłych trajektorii procesów skokowych;
sama statystyczna estymacja należy jednak do tomu IV.
\end{uwaga}

\section{Podsumowanie i połączenie z procesem Wienera}
Proces Poissona miał \emph{stałą} intensywność skoków,
a jego stan zawsze wzrastał o jeden. Łańcuch w czasie
ciągłym dopuszcza zależną od stanu intensywność $q_i$,
różne kierunki zmian oraz pochłanianie. Procesy te nadal
mają trajektorie schodkowe.

W następnych rozdziałach zajmiemy się procesami gaussowskimi
i procesem Wienera, których ciągłe trajektorie różnią się
zasadniczo od trajektorii Poissona. Pomimo tej różnicy
oba światy połączą wspólne narzędzia: filtracje, martyngały,
generatory i równania Kołmogorowa. Jeszcze później
zbudujemy całkę stochastyczną i miarę Wienera na przestrzeni
trajektorii, a następnie zbadamy jej związek z całkami
funkcyjnymi w mechanice kwantowej.


% ==============================================================
% TOM II — ROZDZIAŁ 6: PROCESY GAUSSOWSKIE
% ==============================================================
